\lecture{14}{Mon 06 Oct 2025 12:21}{More submanifolds}

Recall the main theorem. The slice charts theorem tells us that an embedded submanifold has nice charts which look like $(x^1, \ldots ,x^k)\mapsto (x^1, \ldots ,x^k,0, \ldots ,0)$.

\begin{definition}\label{2.1.11}
	Let $p\in M$ and $\Phi : M \to N$ smooth. $p$ is a \emph{regular point} of $\Phi $ if $\mathrm{d}_p\Phi : T_pM \to T_{\Phi (p)} N$ is surjective. Otherwise, $p$ is a \emph{critical point} of $\Phi $. A $c\in N$ is a \emph{regular value} of $\Phi $ if $p$ is a regular point for all $p\in \Phi ^{-1} (c)$. If $c$ is a regular value of $\Phi $, then the fiber of $c$ is a \emph{regular level set} of $\Phi $.
\end{definition}

\begin{corollary}\label{2.1.12}
	Every regular level set between smooth manifolds is a proper embedded submanifold of codimension equal to the dimension of the domain.
\end{corollary}


